\begin{claim}{E1.1}{Proved}{Stability of the pencil for the standard family}
Every sufficiently small deformation of the standard $X_c$ carries a
holomorphic map to $\Pone$ with general fibre a complex two-torus. Its
complete anticanonical system $|-2K|$ is basepoint-free of dimension one
and recovers that map, and its algebraic dimension is one.
\end{claim}
\begin{proof}
D3.1 identifies the semiuniversal germ with the actual $c$-parameter
germ. Thus every sufficiently small deformation is isomorphic to a
member $X_{c'}$ with $\operatorname{Im}c'<-M$. The bundle identity
$K_X=\bar f^*\OO_{\mathcal B}(-6)=\OO_X(-2S_2)$ in D1.2 holds
relatively in this family: all root divisors and every determinant
transition used in its proof vary holomorphically, and their characters
are constant. Consequently $-2K=\bar f^*\OO(12)=f^*\OO_B(1)$
on every member, with a relative line-bundle isomorphism. Connected
fibres imply $f_*\OO_X=\OO_B$, so the two base sections pull back to
the entire two-dimensional section space. They have no common zero.
The resulting map is the constructed $f$ followed by a projective
coordinate change. Its general fibre is a torus by the same local
quotient construction.

We spell out the algebraic-dimension reduction used here. On a
manifold with $H^2(X,\R)=0$, a divisor restricts to a topologically
trivial line bundle on any regular torus fibre. A nonzero effective
divisor on a compact torus has positive degree against a Kahler form,
whereas a topologically trivial class has degree zero. Hence an
effective divisor on $X$ cannot meet a regular fibre in a proper
nonempty divisor: its restriction is either the whole fibre or empty.
The zero and pole divisors of a meromorphic function are therefore
vertical. On a regular fibre where neither contains the fibre, the
restricted function has no pole and is constant. The function hence
descends meromorphically to the base on a dense open set, and the
proper connected-fibre map extends the descent across the finitely
many special values. Thus the meromorphic field is $\C(t)$ and
$a(X)=1$. This argument applies to each small member and does not
appeal to an invalid vanishing theorem.

The conclusion is stronger than what semicontinuity alone gives.
For the original member, projection formula gives
$h^\bullet(f^*\OO(1))=(2,3,1,0)$; the nonzero degree-one group means
that a generic line-bundle extension argument would not establish
stability of its two sections. Here stability follows from the
computed semiuniversal family and its relative root divisors.
\end{proof}

